\documentclass[../../main.tex]{subfiles}
\begin{document}
\begin{oframed}
  \subfile{problem}
\end{oframed}

{\bf [解答]}

全ての自然数は$3$で割りきれるか，$1$または$2$余るかであるから，任意の正の整数$n$の素因数分解を
\begin{align}
  n=3^a\prod_{i=1}^r p_i^{\alpha_i}\prod_{j=1}^s q_j^{\beta_j} \label{eq:1}
\end{align}
と書ける．ただし
\begin{align}
  p_i\equiv1\pmod{3},
  \quad
  q_j\equiv2\pmod{3} \label{eq:2}
\end{align}
であり，要素がない時は積記号は$1$とする．

まず$n$のうちで$3$で割り切れない約数の数$f(n)+g(n)$について考えると，
これは\cref{eq:1}で$p_i,q_j$のみから約数を$A=\prod_{i=1}^r p_i^{\alpha'_i}\prod_{j=1}^s q_j^{\beta'_j}$と構成すれば良い．ただし$\alpha'_i \le \alpha_i, \beta'_j \le \beta_j$である．従ってあり得る組み合わせは各素因数$p_i,q_j$の肩の指数を幾つにするかの組み合わせと等しいから
\begin{align}
  f(n) + g(n) = \prod_i(\alpha_i+1)\prod_j(\beta_j+1) \label{eq:3}
\end{align}
である．

次に$f(n)$について考える．$A$を法$3$の合同式で考えると\cref{eq:2}より
\begin{align}
  A & = \prod_{i=1}^r p_i^{\alpha'_i}\prod_{j=1}^s q_j^{\beta'_j}
  \equiv \prod_{i=1}^r 1^{\alpha'_i} \prod_{j=1}^s (-1)^{\beta'_j}
  \equiv \prod_{j=1}^s (-1)^{\beta'_j} \pmod{3}
\end{align}
だから，これが$A\equiv 1$を満たす時，
\begin{align}
  \beta'_1 + \beta'_2 + \cdots + \beta'_s \equiv 0 \pmod{2}  \label{eq:4}
\end{align}
であれば良い．従ってこのような約数の数は\cref{eq:4}を満たす組み合わせの数を$P$として
\begin{align}
  f(n) = P\prod_i(\alpha_i+1) \label{eq:5}
\end{align}
と書ける．$P$は\cref{eq:4}より，奇数の$\beta'_j$の数が偶数個である場合の数に他ならない．

\paragraph{(1)}

$n=2800=2^4\cdot5^2\cdot7$であるから，\cref{eq:3}は
\begin{align}
  f(n) + g(n) = (4+1)(2+1)(1+1) = 30 \label{eq:6}
\end{align}
である．次に\cref{eq:4,eq:5}について，$2,5$の指数を$\beta'_1,\beta'_2$とすると
\begin{align}
  0 \le \beta'_1 \le 4,
  \quad
  0 \le \beta'_2 \le 2
\end{align}
だから\cref{eq:4}を満たす組み合わせは
\begin{itemize}
  \item $\beta'_1,\beta'_2$ともに偶数の時，$3\cdot2=6$通り
  \item $\beta'_1,\beta'_2$ともに奇数の時，$2\cdot1=2$通り
\end{itemize}
で合計$6+2=8$通りだから\cref{eq:5}より
\begin{align}
  f(n) = 8(1+1) = 16
\end{align}
である．\cref{eq:6}より$g(n)=14$が従うから
\begin{align}
  f(n) = 16,
  \quad
  g(n) = 14
\end{align}
である．

\paragraph{(2)}

\cref{eq:4,eq:5}と同様に考えると，$g(n)$は
\begin{align}
  \beta'_1 + \beta'_2 + \cdots + \beta'_s \equiv 1 \pmod{2}  \label{eq:7}
\end{align}
を満たす組み合わせの数を$Q$として
\begin{align}
  g(n) = Q\prod_i(\alpha_i+1) \label{eq:8}
\end{align}
と表せる．従って$f(n) \ge g(n)$を示すには$P \ge Q$を示せば良い．
ここで\cref{eq:4,eq:7}において和を$j$項目で打ち切った場合に対応する場合の数を$P_j,Q_j$とする：
\begin{align}
  \begin{dcases}
    P_j: \beta'_1 + \beta'_2 + \cdots + \beta'_j \equiv 0 \pmod{2} \\
    Q_j: \beta'_1 + \beta'_2 + \cdots + \beta'_j \equiv 1 \pmod{2}
  \end{dcases}
\end{align}
$P_j =Q_j$または$P_j = Q_j+1$が成立することを$j$に関する数学的帰納法で示す．
そもそも$s=0$の時は$P=1,Q=0$で自明に成立するので以下$s \ge 1$とする．
まず$j=1$の時は$\beta_1$の偶奇に応じて（$k\in\mathbb{Z}\ge0$として）
\begin{align}
  \begin{dcases}
    P_1 = k+1, Q_1 = k   & (\beta_1 = 2k) \\
    P_1 = k+1, Q_1 = k+1 & (\beta_1 = 2k+1)
  \end{dcases} \label{eq:9}
\end{align}
より成立する．そこで$j$での$P_j = Q_j, Q_j+1$の成立を仮定して$j+1$の時を考えると
\begin{align}
  &
  \begin{dcases}
    P_{j+1} = (k+1)P_{j}+kQ_{j} \\
    Q_{j+1} = kP_{j}+(k+1)Q_{j}
  \end{dcases} \quad (\beta_{j+1}=2k) \\
  &
  \begin{dcases}
    P_{j+1} = (k+1)P_{j}+(k+1)Q_{j} \\
    Q_{j+1} = (k+1)P_{j}+(k+1)Q_{j}
  \end{dcases} \quad (\beta_{j+1}=2k+1)
\end{align}
より$j+1$でも題意は成立する．以上より数学的帰納法により$P_j=Q_j, Q_j+1$であり，$j=s$とすれば$P \ge Q$すなわち$f(n) \ge g(n)$であることが示された．\qed

\paragraph{(3)}

\cref{eq:5,eq:8}より
\begin{align}
  \begin{dcases}
    f(n) = P\prod_i(\alpha_i+1) \\
    g(n) = Q\prod_i(\alpha_i+1)
  \end{dcases}
\end{align}
である．よって$g(n)=15=3\cdot5$の時，$Q$がこの約数だから，$Q=1,3,5,15$のいずれかである．(2)より，$P=Q$または$P=Q+1$だから，あり得る組み合わせは以下の通りである．
\begin{align}
  \left(Q,P,\prod_i(\alpha_i+1)\right)
  = (1,1,15), (1,2,15), (3,3,5), (3,4,5), (5,5,3), (5,6,3), (15,15,1), (15,16,1)
\end{align}
このうち$f(n)=PR$が異なる組み合わせとして例えば以下を選択できる．
\begin{align}
  \left(Q,P,\prod_i(\alpha_i+1)\right)
  = (1,1,15), (1,2,15), (3,4,5), (5,6,3), (15,16,1)
\end{align}
それぞれ例えば$n$を以下のように選べば実現できる．以下簡単のため$R=\prod_i(\alpha_i+1)$とおく：
\begin{itemize}
  \item $(Q,P,R)=(1,1,15)$の時，$n=5\cdot 7^{14}$, $f(n)=PR=15$
  \item $(Q,P,R)=(1,2,15)$の時，$n=5^2\cdot 7^{14}$, $f(n)=PR=30$
  \item $(Q,P,R)=(3,4,5)$の時，$n=5^6\cdot 7^4$, $f(n)=PR=20$
  \item $(Q,P,R)=(5,6,3)$の時，$n=5^{10}\cdot 7^2$, $f(n)=PR=18$
  \item $(Q,P,R)=(15,16,1)$の時，$n=5^{30}$, $f(n)=PR=16$
\end{itemize}
従ってとりうる$f(n)$の値は
\begin{align}
  f(n) = 15,16,18,20,30
\end{align}
である．

% TODO: 別解
% $A=\prod_i(\alpha_i+1)$，$T=\prod_j(\beta_j+1)$とおく．3で割り切れない約数について
% \begin{align}
%   f(n)+g(n)=AT,\qquad
%   f(n)-g(n)=A\prod_{j=1}^s\sum_{h=0}^{\beta_j}(-1)^h
% \end{align}
% である．後者は，すべての$\beta_j$が偶数なら$A$であり，奇数の$\beta_j$があれば$0$である
% 奇数の$\beta_j$があれば$f(n)=g(n)=15$となる．すべての$\beta_j$が偶数なら
% \begin{align}
%   15=g(n)=\dfrac{A(T-1)}{2},\qquad f(n)=15+A
% \end{align}
% である．よって$A$は$30$の約数であり，$T$が奇数であることから$A$は$15$の約数である．$A=1,3,5,15$に応じて$T=31,11,7,3$となる．いずれも，$q\equiv2\pmod3$である素数の指数をそれぞれ$30,10,6,2$として実現でき，$A$も$p\equiv1\pmod3$である素数の指数で実現できる．したがって
% \begin{align}
%   f(n)\in\{15,16,18,20,30\}
% \end{align}
% である．\qed
% これは多項式 (1+x+…+x^b) の積に x=−1 を入れると f−g が一発で出る、という母関数を使った見方

\end{document}
